\section{高维数据的低维结构检验}
\label{sec:14-Deep-Learning/07-Interpretability-Mathematics/03}

``数据位于低维流形附近''是深度学习文献中出现频率最高的假设之一，经常被用来解释为什么高维学习居然可行。但这句话本身需要一个操作化的检验方案，否则就变成了训练成功后的事后解释——学好了是因为低维，低维的证据是学好了。更隐蔽的混淆是：随机投影能够降维并保持距离，这一事实本身不能论证数据具有低内在维度。

\subsection{JL 引理不要求流形假设}

\begin{theorem}[Johnson--Lindenstrauss 引理的推论]
对欧氏空间中的 \(n\) 个点，存在到

\[
k=O\!\left(\varepsilon^{-2}\log n\right)
\]

维空间的线性映射，使所有点对距离都保持在原距离的 \(1\pm\varepsilon\) 倍内。随机 Gaussian 投影或适当稀疏投影以高概率实现这一点。
\end{theorem}

JL 的目标维数只依赖点数和允许的失真度，不依赖原始维数，也不假设这些点来自某个低维流形。它说的是：任意有限距离表都可以被近似保存进对数级压缩的维度中。这个结论对理解随机投影很有用，但它不识别生成数据的潜变量，不保证新样本、概率密度或拓扑结构同样被保留。

因此，拿随机投影把数据压到 \(k\) 维后模型表现不变，只能说明当前任务对这次压缩还算稳健。它不是``数据的内在维度等于 \(k\)''的证明。

\subsection{流形假设需要局部和尺度信息}

如果数据在某个尺度上确实接近 \(d_{\mathrm{int}}\) 维光滑流形，那么在足够小的邻域里，变化应主要落在约 \(d_{\mathrm{int}}\) 个切向方向上。常用的诊断工具包括：

\begin{itemize}
\tightlist
\item 局部 PCA：取每个点的邻域，对邻域协方差做谱分解，观察特征值是否在某个位置出现稳定陡降；
\item 最近邻尺度律：邻域半径随样本量变化的速率应与某个维数模型的预测一致；
\item 相关维数或参与率：从距离计数或谱集中程度构造一个有效维数的标量摘要；
\item 多尺度重复：改变邻域大小、距离度量与预处理方式后，检查维数估计是否保持稳定。
\end{itemize}

这些方法估计的都是特定尺度、特定度量和特定噪声模型下的有效维数。邻域取得太小，被采样稀疏性支配——局部 PCA 看起来像满维；取得太大，又把弯曲、扭结或不同分支的数据混在一起碾平。两个极端之间需要一个平稳区间，其存在本身就是证据的一部分。

\subsection{一个最小的反例检查}

设想信号位于二维平面，但每个坐标都叠加了小幅独立高维噪声。在极小的尺度上，最近邻差异主要由噪声主导，局部 PCA 的谱看起来接近满维；在较大的尺度上，两个信号方向才逐渐浮出占据主导。问``数据到底几维''没有一个脱离尺度的唯一答案。

反过来，一条在空间中缠绕得非常复杂的一维曲线，虽然内在维数为 1，仍可能需要大量样本才能覆盖其几何，也可能具有高度振荡的分类边界。低维不自动等于容易学习。

\subsection{任务维数与数据维数是两回事}

预测目标可能只依赖数据变化中的少数方向，即使完整数据分布的支撑并不低维。可以通过充分降维、受控特征删除、跨环境预测稳定性或表示干预，检查哪些方向对当前任务必要。只看无监督重构方差，可能回收了高变异但与任务毫无关系的因素。

深度网络也不会因为架构深就自动发现真实低维结构。卷积写入了局部参数共享，注意力写入了内容相关的读取模式；训练在这些结构约束下搜索有效的表示。但最终得到的表示是否低维、可迁移、且符合数据生成机制，需要独立诊断而不是默认相信。

\subsection{可以相信什么}

如果多种维数估计在合理的尺度范围内稳定，低维模型在留出数据上保持了距离、重构或预测质量，并且结论对度量选择和噪声处理不敏感，你就获得了支持低维近似的证据。即便如此，仍应把结论限定在当前数据范围、当前尺度与当前任务内，而不是宣称世界存在一个唯一的低维真相。

下一篇讨论模型归因时会遇到同一种原则：先精确定义你要解释的对象，再说明你的估计对哪些扰动稳定；可视化和单一分数只能提供线索，它们本身不是检验。
